\documentclass[11pt]{article}

\usepackage{amsmath,amssymb}
\usepackage{xurl}
\usepackage[hidelinks]{hyperref}
\setlength{\emergencystretch}{2em}

% Shared commands for notes in docs/paper-gaps/.
% Notation follows the LDT paper (references/ldt-paper/).

\newcommand{\Lean}{\textsc{Lean}}
\newcommand{\leanid}[1]{\textsf{#1}}
\newcommand{\ghissue}[1]{\href{https://github.com/LionSR/MIPStarRE/issues/#1}{GitHub issue \##1}}
\newcommand{\ghpr}[1]{\href{https://github.com/LionSR/MIPStarRE/pull/#1}{GitHub PR \##1}}

% --- Paper-compatible notation ---
\newcommand{\F}{\mathbb{F}}
\newcommand{\N}{\mathbb{N}}
\newcommand{\C}{\mathbb{C}}
\newcommand{\tr}{\operatorname{tr}}
\newcommand{\ot}{\otimes}
\newcommand{\E}{\mathop{\bf E\/}}
\newcommand{\bx}{\mathbf{x}}
\newcommand{\by}{\mathbf{y}}
\newcommand{\bu}{\mathbf{u}}
\newcommand{\bv}{\mathbf{v}}

% Approximation relation (paper uses \approx_\eps and \simeq_\zeta)
\newcommand{\approxeps}{\approx_{\varepsilon}}
\newcommand{\simgeq}{\simeq_{\zeta}}

% Polynomial spaces (paper uses \polyfunc, \polysub, \polymeas)
\newcommand{\polyfunc}[3]{\operatorname{Poly}(#1,#2,#3)}
\newcommand{\polysub}[3]{\operatorname{PolySub}(#1,#2,#3)}
\newcommand{\polymeas}[3]{\operatorname{PolyMeas}(#1,#2,#3)}


\title{The Large-\texorpdfstring{$k$}{k} Side Condition in the Main Induction}
\author{}
\date{2026-04-29}

\begin{document}
\maketitle

\section{The assertion}

This note concerns the \(k\)-hypothesis in the main induction theorem of
\cite[Section~\emph{The main induction step}]{Ji2020LowIndividualDegree}, and
its use in the proof of the main formal soundness theorem.
\footnote{The original discrepancy was recorded in
\href{https://github.com/LionSR/MIPStarRE/issues/906}{GitHub issue \#906} and
implemented in
\href{https://github.com/LionSR/MIPStarRE/pull/912}{GitHub pull request \#912}.
The documentation issue for this note is
\href{https://github.com/LionSR/MIPStarRE/issues/928}{GitHub issue \#928}.  In
the TeX source consulted here, the main induction theorem is in
\path{references/ldt-paper/inductive_step.tex}, lines 7--18; the successor step
is in the same file, lines 429--568; and the pasting theorem used there is in
lines 299--330.}
Here \(m\) is the number of variables, \(q\) is the field size, \(d\) is the
individual degree bound, and \(k\) is the sampling parameter in the consistency
estimates.  The integer \(k\) is carried unchanged through the induction on
\(m\).

In the theorem as stated, an \((\varepsilon,\delta,\gamma)\)-good symmetric
strategy for the \((m,q,d)\) low individual degree test is assumed to satisfy
only
\begin{equation*}
  k\ge md.
  \tag{\path{thm:main-induction}}
\end{equation*}
Under this hypothesis, the theorem produces a polynomial measurement
\(G\in\operatorname{PolyMeas}(m,q,d)\) satisfying the advertised consistency
estimate.  The same \(k\)-condition is then inherited by the stated main formal
soundness theorem after the reduction from a general two-prover strategy to a
symmetric one.

\section{The use of pasting}

The successor step proves the induction theorem in dimension \(m+1\) assuming
the theorem in dimension \(m\).  In that step the ambient strategy is a strategy
for the \((m+1,q,d)\) test, and the proof assumes
\begin{equation*}
  k\ge (m+1)d.
  \tag{\path{sec:proof-of-main-induction}}
\end{equation*}
This is sufficient to apply the induction hypothesis to each \(x\)-restricted
strategy, since the restricted strategies have dimension \(m\), and
\((m+1)d\ge md\).

After the restricted measurements have been improved to projective
submeasurements, the proof pastes them into a measurement in dimension \(m+1\).
The pasting theorem invoked at this point has its own size hypothesis.  In the
notation of the successor step, it requires
\begin{equation*}
  k\ge 400md.
  \tag{\path{thm:ld-pasting-in-induction-section}}
\end{equation*}
This condition is not a consequence of \(k\ge (m+1)d\).  For instance, when
\(m=1\) and \(d>0\), the latter allows \(k=2d\), whereas the pasting theorem
requires \(k\ge 400d\).  Thus the argument as written does not justify the
successor step throughout the interval in which \(k\ge (m+1)d\) but
\(k<400md\).

Equivalently, the missing implication is
\[
  k\ge (m+1)d
  \quad\Longrightarrow\quad
  k\ge 400md,
\]
which is false for the relevant range of parameters.  The issue is therefore a
missing side condition, not a notational convention or a peculiarity of the
formal language.

\section{Current formalization status}

The current blueprint makes the paper statement visible again: the displayed
main induction theorem is stated with \(k\ge md\), and the displayed successor
step assumes \(k\ge (m+1)d\).  This is the statement in the TeX source of the
paper, but the blueprint does not mark it as proved in Lean.  No present Lean
declaration has this theorem as its public type.

The checked successor-step assembly is instead recorded as a separate
Lean-only conditional proposition.  That proposition retains the side
conditions actually used by the current proof route, including \(1\le k\) and
\(400md\le k\).  Downstream handoff statements used for the main formal
assembly also retain the corresponding large-\(k\) hypothesis.  These
conditional statements are not proofs of the paper theorem in the range
\(md\le k<400md\); they expose the extra side condition as an open proof
obligation rather than incorporating it into the paper statement.
\footnote{The relevant blueprint statements are
\path{blueprint/src/chapter/ch02_test.tex}, lines 107--140, and
\path{blueprint/src/chapter/ch10_induction.tex}, lines 197--245 and
346--354.  The Lean conditional successor-step wrapper is
\path{MIPStarRE.LDT.MainInductionStep.mainInductionPublicWrapper}; the pasting
theorem is \path{MIPStarRE.LDT.Pasting.Bernoulli.ldPasting}, with the
corresponding large-\(k\) hypothesis.}

\section{Conclusion}

The verdict is that
\href{https://github.com/LionSR/MIPStarRE/issues/906}{issue \#906} records a
genuine theorem-level side-condition gap in the stated proof.  The hypothesis
\(k\ge md\), or in the successor step \(k\ge (m+1)d\), does not imply the
large-\(k\) hypothesis required by the pasting theorem.  Pull request
\href{https://github.com/LionSR/MIPStarRE/pull/912}{\#912} repaired the formal
route by strengthening the displayed hypothesis.  The present blueprint policy
instead keeps the paper theorem statement visible and records the stronger
condition on the Lean-only conditional wrapper and on the downstream consumers
that still use the current pasting route.

This conclusion does not assert that the original theorem is false for
\(md\le k<400md\).  It says that the proof route through the stated pasting
theorem does not establish that range.  A proof of the smaller-\(k\) range would
require an additional argument, such as a genuine saturation lemma for that
interval or a pasting theorem with a correspondingly weaker size hypothesis.

\bibliographystyle{alpha}
\bibliography{references}

\end{document}
